\section{Actual three/four arms with unrestricted four count once $r\ge50$}
This theorem removes the preceding ternary-surplus restriction.
It remains a specified actual mixed-capacity family; the GENERAL
one-descent conjecture is not asserted.

\begin{theorem}
For EVERY $\alpha\ge1$, EVERY $r\ge50$ and EVERY $s\ge0$,
the actual arm $1^\alpha3^r4^s$ has complete C at every current
rank and every shift. There is NO ratio constraint between $r$ and $s$.
More precisely, writing $s_0=\lceil(r+2)/3\rceil$:
\begin{itemize}
\item if $s\ge s_0$, the ENTIRE ordinary coefficient tail is LC
at every centre $N\ge s+4$;
\item whenever $r+4\ge3s$, the ENTIRE ordinary coefficient tail is
LC at every centre $N\ge r+s+2$.
\end{itemize}
Both counts and padding are unbounded.
\end{theorem}

\subsection{A stronger near-boundary payment, valid at every relevant rank}
Keep $b=1/\lambda_3$, $1/2\le b\le6/11$.
For EVERY $S\ge18$ and EVERY $n$ such that $3S\le n+2$,
\[
                         10\,2^S b^{\lfloor n/2\rfloor}<1.       \tag{1}
\]
Indeed, with $k=\lfloor n/2\rfloor$, $2k+3\ge3S$, so
\[
 (10\,2^Sb^k)^2\le800(4b^3)^S
 \le800(2/3)^{18}<1.
\]
The factor eight uses $b^3\ge1/8$. The whole natural-index/floor/power
payment is kernel checked; no limiting argument replaces any rank.

Fix $r\ge50$ and $r+4\ge3s$, and put $S=\max(s,18)$.
Then $r+4\ge3S$, $S\ge18$, $r\ge32$.
The already proved proper-quotient cone at count $S$ puts
$b^{r+1}4^S\le1/64$, and monotonicity gives the same for the ACTUAL
smaller count $s$. Thus the actual proper quotient $Q$ remains positive
globally LC with both remainders in the quarter box.
For $n=r+2$, equation (1) applies. Monotonicity of $2^s$ transfers
it to the actual count. The preceding exact cubic norm therefore gives
the scaled ordinary-minor bound and strict root-gap bound required
in the previous actual-cut proof.

We spell out why that proof's other hypotheses survive this refinement.
The corrected cumulative $V=Q/(1-z)+\gamma/(1-z)$, its positive LC
primitive after rank zero, its cumulative LC seed, the TRUE degree-one
source $C_E=BV+(R_0-\gamma)+b\gamma z$, all three finite-padding
startups, bounded future inverse and actual mode guard are unchanged.
They require only the quarter box, $w=r+s\ge32$ and
$\deg K\le A+1$, with $A=w+2$. The last condition is exactly
$r+4\ge3s$, as shown by the earlier complete degree calculation.

Only the physical first-boundary minor needs a further distinction.
If $r+4=3s$, then $A=4s-2$ is the LAST centre in the NEW reference's
proved ordinary head. Hence $M_{R_s}(A)\ge0$ directly.
Otherwise $r+3\ge3s$, and all three reference reads are purely
$2^s h_3$; (1) gives $M_{R_s}(A)\ge-(R_s)_{A-1}$.
In either case the preceding first-actual-defect kernel applies to
the TRUE loss one or two and its TRUE next loss. No extrapolation
of the future model to offset $-1$ is used.
Every step of the signed-cut proof therefore proves the claimed
ordinary tail from $A$ for all $r\ge50$, $r+4\ge3s$.

If the stronger $r+3\ge3s$ holds, the two reference entry reads are
homogeneous as well. Equation (1) and the exact ternary residual norm
give the strict ACTUAL entry $F_A<\lambda_3F_{A-1}$.
The proven complete reference C, its positive root-cumulative gaps
and the bounded ordinary-to-C bridge then give FULL actual C.
Thus we already have, at ALL counts,
\[
 r\ge50,\quad r+3\ge3s\quad\Longrightarrow\quad\text{complete C}. \tag{2}
\]
This refinement uses the previous complete actual cut assembly;
it does not infer actual C merely from ordinary LC of a polynomial
quotient.

\subsection{An analytic unbounded diagonal seed}
For arbitrary fixed $r\ge50$, let
\[
 s_0=\left\lceil{r+2\over3}\right\rceil
                     =\left\lfloor{r+4\over3}\right\rfloor\ge18.
\]
Then $3s_0-4\le r\le3s_0-2$. The preceding near-boundary tail
theorem proves LC from $A_3=r+s_0+2$, for EVERY initial padding.
At this entrance, $A_3$ is respectively $W-2,W-1,W$, with $W=4s_0$.

Before $A_3$ the actual arm agrees exactly with the infinite-three
reference $R_{s_0}$. Its proved ordinary head covers
$s_0+4\le n\le W-2$. Therefore all centres from $s_0+4$ through
$A_3-2$ are covered, and lowering the first NEXT read at $A_3-1$
improves its minor whenever that centre is still in this head.

The ONLY remaining case is $r=3s_0-2$, $A_3=W$, whose preceding
centre is $W-1$. All three reference reads there are homogeneous.
Use (1) at $n=r+1=3s_0-1$ to get its scaled unit minor bound.
Lowering its next coefficient by the compulsory integer loss
$\delta\ge1$ then pays it EXACTLY:
\[
 M_{\rm actual}(A_3-1)
       =M_{R_{s_0}}(A_3-1)+\delta(R_{s_0})_{A_3-2}\ge0.
\]
This last joining edge is a separate scoped kernel.
The entire ordinary actual tail from $s_0+4$ is now proved.
This is an ANALYTIC proof on the unbounded diagonal $s_0(r)$,
not a finite diagonal used as an induction premise.

\subsection{All larger four counts and the complete phase cover}
For $r\ge50$, the finite seed $F_{1^\alpha3^r}$ agrees with $h_3$
through rank seven at EVERY $\alpha\ge1$: its first difference from
$h_3$ is only at $r+2$. Its first four losses relative to $h_4$ are
therefore EXACTLY $(1,2,5,12)$ at ranks $4,5,6,7$.

The existing ordinary four-convolution heredity now applies to this
seed at every count at least $s_0\ge18$. Its entire requirements are
the first four losses, the ordinary $h_4$ unit minor and factor-two
growth bounds, the three established startup budgets, and LC of the
old ordinary tail. It has no premise that the smaller seed alphabet
must be binary. Concretely, the three new left centres are paid by
the generic startup kernel with $(a,b,c,d)=(1,2,5,12)$.
For all later centres truncate the old LC tail below its first
three-read rank and convolve with $[4]$; all new reads agree.
The old tail extends through its zero end, so there are no unsigned
terminal source corrections in this actual full-tail induction.
The new atom at physical zero is outside the whole read band.

Induction proves ordinary coefficient LC from $s+4$ for EVERY
$s\ge s_0$, at every padding. This rank is precisely the actual first
loss from $h_4$, of size one. The proved universal first-defect
ordinary criterion gives full actual C throughout this region.
For $s<s_0$, the integer seed arithmetic gives $3s\le r+1$, so (2)
already proves full C. These two cases cover EVERY $s\ge0$.
This completes the theorem without a new finite proof premise.

\begin{corollary}[Original A7 family, with no three/four count ratio]
Adjoin ALL $1^\alpha3^r4^s$ with $r\ge50$, $s\ge0$ to $\mathcal A_6$,
and call the union $\mathcal A_7$. It is closed under final-part
deletion: removing a four preserves $r$, and once no fours remain,
all successive pure-ternary arms are already in the old class.
Every member has complete C and first missing weight at most five.
The existing sharp first-missing-five theorem therefore proves FULL
pair A2, the entire positive-source cone for both nonrectangular arms,
and ORIGINAL UNAMPLIFIED coefficient unimodality of the strict core
and BOTH relatives at EVERY strict central peak. One arbitrary
legal heavier final append is also allowed.
Within a single arm, the number of fours is now unrestricted relative
to the unbounded number of threes $r\ge50$.
\end{corollary}

\paragraph{Verification scope and remaining work.}
Four new scoped Lean reports check the entire near-boundary damped
payment, analytic diagonal/phase-cover arithmetic, monotonicity of
both norm and tilt with count, and the compulsory next-only loss
join. All complete cut/source/startup/mode/quotient/reference and
first-missing-five kernels are reused without recounting. The series
and LC-convolution assemblies and original-family corollary remain
written proofs, not end-to-end Lean composition formalizations.
There is no new finite entrance premise. Extra actual and greedy
distinct-composition enumerations are QA only.
The finitely many three counts $r<50$ with arbitrary four count,
binary-prefix mixed arms $1^\alpha2^u3^r4^s$, and general capacities
remain outside this theorem. The GENERAL original conjecture is OPEN.
